\documentclass[12pt,a4paper]{article}
\usepackage{ugmath}
\usepackage{float}
\usepackage{placeins}
\usepackage{array}
\usepackage[labelfont=bf,labelsep=space]{caption}
\newcommand{\paren}[1]{\left( #1 \right)}
\newcommand{\alumno}{Ricardo León Martínez}
\newcommand{\materia}{Algebra Lineal I}
\newcommand{\profesor}{Claudia Reynoso Alcántara}
\newcommand{\tarea}{Tarea 8}
\newcommand{\fecha}{23/3/2026}


\begin{document}

    \begin{center}
        {\large\textbf{UNIVERSIDAD DE GUANAJUATO}}\\[0.3cm]
        {\normalsize\textbf{DIVISIÓN DE CIENCIAS NATURALES Y EXACTAS}}\\
        {\normalsize\textbf{CAMPUS GUANAJUATO}}\\[1cm]

        {\Large\textbf{\tarea\ (\materia)}}\\[1cm]
    \end{center}

    \noindent
    \textbf{Nombre:} \alumno 
    \hfill 
    \textbf{Fecha:} \fecha 
    \hfill 
    \textbf{Calificación:} \rule{3cm}{0.4pt}

    \vspace{0.3cm}

    \begin{mdframed}
        Dado un espacio vectorial $V$ de dimensión finita $n$ sobre $F$ y dada una base ordenada $\mathcal{B}$ de $V$. Demuestra
        que para todo $v,w\in V$ y $a\in F$, se tiene que:
        \begin{equation*}
            [av+w]_{\mathcal{B}}=a[v]_{\mathcal{B}}+[w]_{\mathcal{B}}.
        \end{equation*}
    \end{mdframed}

    \begin{mdframed}
        Demuestra que $\mathcal{B}=\{(1,2,0,1),(0,1,1,3),(2,0,1,1),(1,1,1,1)\}$ es una base de $\R^{4}$. Encuentra las coordenadas
        de $(a,b,c)\in\R^{4}$ en la base ordenada $\mathcal{B}$.
    \end{mdframed}

    \begin{mdframed}
        Considera $V_{2}$ el espacio vectorial sobre $\R$ de polinomios de grado menor o igual a 2 con coeficientes en $\R$. Sea
        $c\in\R\setminus\{0\}$, entonces $\BB_{2}=\{1,x+c,x^{2}+c^{2}\}$ es una base ordenada para $W_{2}$. Encuentra la matriz $P$
        de cambio de base de $\BB_{2}$ a la base ordenada $\BB_{1}=\{1,x,x^{2}\}$.
    \end{mdframed}

    \begin{mdframed}
        Demuestra que $\BB_{2}=\{(2i,1,0),(2,-1,1),(0,1+i,1-i)\}$ es una base ordenada de $\C^{3}$ y encuentra la matriz de cambio de base
        de $\BB_{2}$ a la base canonica $\BB_{1}$.
    \end{mdframed}

    \begin{mdframed}
        Sea $V$ un espacio vectorial de dimensión finita $n$ sobre un campo $F$
    \end{mdframed}
\end{document}